% 09_concepts.tex -- self-contained "concept school": every idea the reader
% needs to understand Section~\ref{sec:experiments}, taught from zero with
% this project's own files as the worked examples. Video + paper references
% live in the bibliography at the end of this file.

\section{Concept school: the ideas behind the experiments}
\label{sec:concepts}

This section teaches, from zero, each idea the experimental evaluation
(Section~\ref{sec:experiments}) rests on --- reactive obstacle avoidance,
histogram steering, reasoning under sensor uncertainty, Monte Carlo
experiment grids, and the ToF-vs-LiDAR sensor trade-off. Each subsection ends
with pointers for going deeper: the paper that introduced the idea, a
textbook chapter, and at least one video lecture. The worked examples are
always this project's own files, so every concept can be seen running in
code you already have.

\begin{keybox}{Concept map --- where each idea lives in this project}
\footnotesize
\begin{tabular}{@{}p{2.9cm}p{4.6cm}p{2.6cm}p{3.1cm}@{}}
\toprule
\textbf{Concept} & \textbf{Implemented in} & \textbf{Tested by} &\textbf{Taught here} \\
\midrule
Potential fields (APF) & \code{apf\_force.m} / \code{vfh.py} reflex & E1 (local minimum) & \S\ref{sec:concept-apf} \\
Histogram (VFH) & \code{pvfh\_scan.m} + \code{steering\_dispatch.m} / \code{vfh.py} & E1--E3 baselines & \S\ref{sec:concept-vfh} \\
Sensor noise model & \code{sensor\_noise.m} / \code{vfh.py::sensor\_noise} & E2, E3, E6 & \S\ref{sec:concept-unc} \\
Probabilistic blockage & \code{pvfh\_prob.m} / \code{vfh.py::pvfh\_prob} & E2 (headline) & \S\ref{sec:concept-unc} \\
Monte Carlo grid & \code{run\_paper\_experiments.m} & E1--E14 & \S\ref{sec:concept-mc} \\
Paired seeds & same PRNG seed per run & E2, E5, E7 & \S\ref{sec:concept-mc} \\
ToF vs LiDAR & \code{sensor\_noise.m} \code{'tof'} vs \code{'lidar'} & E3, E6 & \S\ref{sec:concept-tof} \\
Formation geometry & \code{formation\_geometry.m} / \code{geometry.py} & E4, E5 & \S\ref{sec:concept-vfh} \\
\bottomrule
\end{tabular}
\end{keybox}

\subsection{Reactive obstacle avoidance I: potential fields}
\label{sec:concept-apf}

The oldest answer to ``how should a robot steer around things it can see''
is to pretend space is a landscape of hills and a valley. The goal exerts an
\emph{attractive} force, every obstacle an \emph{inflating} repulsion, and
the robot simply walks downhill on the sum:
\begin{equation}
u = k_{\mathrm{att}}\,(x_{\mathrm{goal}} - x)\; -\;
    \sum_{i} k_{\mathrm{rep}}\,\nabla \phi(d_i),
\qquad
\phi(d) \approx \Big(\tfrac{1}{d}-\tfrac{1}{d_0}\Big)^2 \ \text{for } d < d_0,
\label{eq:apf-concept}
\end{equation}
where $d_i$ is the clearance to obstacle $i$ and $d_0$ the influence radius.
Three properties matter for everything that follows:

\begin{enumerate}[leftmargin=1.9em]
\item \textbf{It is reactive.} No map, no plan --- the command is a function
of the \emph{current} scan only. That is why it costs microseconds (E1's
\code{apf} compute bar) and why it can be reflexively layered under any other
planner.
\item \textbf{It has a famous failure mode.} Two obstacle faces pushing from
opposite sides can cancel the goal attraction exactly: a \emph{local
minimum}. The robot then presses into a wall forever. This is not a bug in
one implementation; it is a theorem about eq.~\eqref{eq:apf-concept}, shown
rigorously by Koren and Borenstein \cite{koren1991}. E1's \code{apf} row
(Table~\ref{tab:e1}) is where to look for this theorem in the data: a law
that stops making progress without ever being refused by the safety gate.
\item \textbf{It fixes itself if you add a histogram.} Vector fields answer
``which way do forces point''; histograms (next subsection) answer ``which
bearings are free'', and a free-bearing search cannot be trapped the same
way.
\end{enumerate}

\paragraph{Go deeper.} Khatib introduced the real-time potential-field
method \cite{khatib1986}; Koren and Borenstein catalogued its pathologies
\cite{koren1991}. The friendliest video treatment is Chapter~10.6 of the
Northwestern \emph{Modern Robotics} course \cite{lynchvideo}, six minutes
that build eq.~\eqref{eq:apf-concept} visually; the textbook chapter behind
it is Lynch \& Park \cite{lynch2017}, \S10.6.

\subsection{Reactive obstacle avoidance II: histograms in motion (VFH)}
\label{sec:concept-vfh}

VFH replaces vector addition with \emph{voting}. Take the $n$ range beams of
the scan. Each beam $b$ measures range $d_b$ along bearing $\theta_b$. Reduce
the plane to $\alpha$ sectors of $\Delta\theta$ degrees and give each sector
an obstacle density: the classic choice, used verbatim in
\code{pvfh\_scan.m}/\code{vfh.py}, is
\begin{equation}
h_b \;=\; \Big(1-\tfrac{\min(d_b, d_{\max})}{d_{\max}}\Big)^2,
\qquad
h_k = \sum_{b \in \mathrm{sector}\,k} h_b,
\label{eq:vfh-concept}
\end{equation}
so a close return votes strongly, a far return weakly, and each sector sums
its beams' votes. Sector $k$ is \emph{blocked} when $h_k$ exceeds the
threshold \code{tau}; the robot steers into the \emph{gapped} direction
closest to the goal bearing that also respects the robot's turning radius
(a ``mask'' on candidate directions). The pipeline is exactly:
scan $\to$ per-beam kernel \eqref{eq:vfh-concept} $\to$ sector histogram
$\to$ threshold $\to$ masked gap selection $\to$ heading setpoint.

Why this cannot share APF's failure mode: the two opposing walls both vote
``blocked'' for their own sectors, but the \emph{search} over free sectors
sees the aisle between them and takes it. The failure changes character: VFH
cannot get stuck, but it \emph{can} misread --- which is where noise enters.

\paragraph{How to read any VFH plot.} A polar histogram is a wheel with one
wedge per sector; wedge height $=$ density \eqref{eq:vfh-concept}, the
horizontal line is the threshold \code{tau}, and the chosen heading is the
free wedge nearest the goal direction (drawn as an arrow). When you run
\code{Simulation8\_dynamic.m} interactively with \code{viz.enable=true}, this
wheel is the live figure to watch; the PDF flowchart in
Section~\ref{sec:pvfhflow} shows the same steps statically.

\paragraph{Go deeper.} The original paper \cite{borenstein1991} is short,
heavily illustrated and still the best read --- free PDF from Yoram Koren's
Michigan page. \emph{VFH+} \cite{ulrich1998} adds the robot-radius mask and
hysteresis our implementation inherits. The MathWorks documentation page
\cite{mwvfh} shows MATLAB's own \code{controlVFH+} object, useful to compare
parameter names. For video: a lecture on sonar sensing and VFH
\cite{vfhvideo} and the shorter ``Vector Field Histogram'' explainer
\cite{vfhvideo2}; the Wikipedia article \cite{vfhwiki} is a good 5-minute
recap with the original figures.

\subsection{Reasoning under sensor uncertainty: the probabilistic histogram}
\label{sec:concept-unc}

Deterministic VFH answers one question per beam: \emph{is the bearing
free?} --- threshold \eqref{eq:vfh-concept} against \code{tau}. With a
perfect sensor that is fine. With real sensors it fails twice:

\begin{enumerate}[leftmargin=1.9em]
\item \textbf{Missed obstacle (dangerous).} A wall at 1.6~m whose return
lands \emph{inside} the reaction envelope can still read ``free'' if the
measurement noise pushes the estimate out --- the robot drives at a wall it
just saw. This is the regime E2 probes: the high-$\sigma_0$ columns of
Table~\ref{tab:e2}, where the gate-rejection count measures exactly these
``drove at something it had just seen'' events.
\item \textbf{Flicker (wasteful).} Graded kernel values near \code{tau}
flip between free and blocked as noise arrives, so the robot brakes and
lurches --- clearance collapses even when it does not collide.
\end{enumerate}

The probabilistic layer answers a better question: \emph{given everything I
know --- range estimate, sensor variance, FoV geometry --- how likely is this
bearing to be blocked?} Three ingredients, one per answer:

\paragraph{1. A per-beam measurement distribution.} Each beam returns
$\hat d_b$ modelled as true range plus Gaussian noise, $\hat d_b \sim
\mathcal N(d_b,\, s_b^2)$, where \code{sensor\_noise.m} sets $s_b$: a flat
$\sigma_0$ for the lidar kind, and for the ToF kind a linear growth term
$s_b = \sigma_0 + k_{\mathrm{rng}}\, d_b$ with a $60^\circ$ FoV arc ---
matched to how cheap ToF arrays actually behave (\S\ref{sec:concept-tof}).

\paragraph{2. A blockage probability.} The bearing is ``blocking'' if the
\emph{true} obstacle surface lies inside the reaction envelope, i.e.\ true
range $d_b \le d_{\mathrm{block}}$. Since $\hat d_b$ is a noisy proxy for
$d_b$, that event has probability
\begin{equation}
p_b \;=\; \Pr(d_b \le d_{\mathrm{block}} \mid \hat d_b, s_b)
       \;=\; \Phi\!\Big(\frac{d_{\mathrm{block}} - \hat d_b}{s_b}\Big),
\label{eq:prise-concept}
\end{equation}
with $\Phi$ the standard normal CDF. If \eqref{eq:prise-concept} looks
familiar, it is exactly the beam model of probabilistic robotics
\cite{thrun2005} --- the same $z$-test against a range threshold, turned
into a soft membership. The floor $p_{\mathrm{unknown}} = 0.05$ applies
outside the sensor's FoV: the unobserved rear hemisphere is neither a phantom
wall nor a blind spot, it is a calibrated shrug.

\paragraph{3. Sector pooling and the expected kernel.} A sector is as
blocked as its \emph{best evidence}: independent per-beam probabilities pool
by noisy-OR,
\begin{equation}
p_{\mathrm{sector}} \;=\; 1 - \prod_{b \in k} (1 - p_b),
\label{eq:noisyor-concept}
\end{equation}
so every additional beam covering the sector genuinely adds evidence. In
the implementation this pooled probability is a \emph{diagnostic} (it is
logged and plotted, and it is the quantity a risk monitor would threshold);
the steering itself uses the kernel below, for the saturation reason given
in \S\ref{sec:pvfh}. Finally the deterministic kernel \eqref{eq:vfh-concept} is replaced by its
\emph{posterior expectation},
\begin{equation}
u_b \;=\; \sqrt{\tfrac{\sigma_h^2}{\sigma_h^2 + s_b^2}}\;
        \exp\!\Big(\!-\tfrac{d^2}{2(\sigma_h^2 + s_b^2)}\Big),
\label{eq:ukernel-concept}
\end{equation}
i.e.\ the expected value of the Gaussian danger kernel when the range is a
belief rather than a fact. Equation \eqref{eq:ukernel-concept} has two
properties that matter more than its formula:

\begin{itemize}[leftmargin=1.6em]
\item \textbf{Exact deterministic limit.} As $s_b \to 0$ the prefactor
$\to 1$ and \eqref{eq:ukernel-concept} $\to$ \eqref{eq:vfh-concept}: with a
perfect sensor P-VFH \emph{is} the deterministic histogram
(Proposition~1). The paired-difference plot of
Section~\ref{sec:paired} is where this shows up as data: the P-VFH~$-$~pVFH
difference should shrink toward zero at the low-noise end.
\item \textbf{A reaction envelope that widens with noise, never
saturation.} The range at which a return starts to block grows with $s_b$
(Proposition~2, Fig.~\ref{fig:theory}) --- uncertainty makes
the robot react \emph{earlier}. Note the precise statement: the kernel
\emph{value} rises with noise only for returns beyond
$\sqrt{\sigma_h^2 + s_b^2}$; very near returns read slightly lower but stay
far above the threshold. (An earlier version of this report claimed the
kernel inflates at every range; that was wrong.)
\end{itemize}

The one-sentence summary: \textbf{VFH thresholds a histogram of measurements;
P-VFH thresholds a histogram of beliefs.}

\paragraph{Go deeper.} Thrun, Burgard \& Fox, \emph{Probabilistic Robotics}
\cite{thrun2005}: Chapter~2 (recursive state estimation) and the beam
models of Chapter~6 are the mathematical home of \eqref{eq:prise-concept}.
Cyrill Stachniss's video lectures are the standard free companion: the Bayes
filter lecture \cite{stachnissbayes} builds the inference behind ``probability
of a hypothesis given a noisy measurement'', and his 5-minute occupancy-grid
explainer \cite{stachnissgrid} shows the identical free/blocked/unknown
triage P-VFH applies per sector.

\subsection{Monte Carlo experiments: factors, grids, cells, seeds}
\label{sec:concept-mc}

A robot run is stochastic --- noise realisations differ run to run --- so one
run proves nothing. The fix is to run the mission many times under a
controlled experimental grid and report statistics. Section~\ref{sec:gridconcept}
gives the operational vocabulary; here is the reasoning behind it.

\begin{keybox}{The five words, in one breath}
A \textbf{factor} is what you vary on purpose; its values are
\textbf{levels}; the \textbf{grid} is every combination of levels; one
combination is a \textbf{cell}; one full mission executed in a cell under
one \textbf{seed} (a pinned random number generator) is a \textbf{run}.
Repeat runs with different seeds and report means/rates/spreads:
that repetition is the \textbf{Monte Carlo} method.
\end{keybox}

\paragraph{Why the seed is the most important knob.} Seeding the PRNG
(``random'' numbers drawn from a reproducible generator) makes a stochastic
experiment \emph{replayable}. Every run is seeded explicitly (MATLAB \code{rng(seed)}; Octave
\code{randn('seed',...)}; the Python port \code{np.random.default\_rng(seed)}). Consequences you can rely on:
re-running a cell reproduces its CSV byte-for-byte (a regression test in
itself), and --- the design trick --- \emph{two algorithms sharing a seed
see the identical noise realisation}. Their per-seed difference then
contains no noise variance at all, only the steering-law difference: a
paired comparison, the experimental analogue of a paired t-test. That is
why Section~\ref{sec:paired} can resolve small clearance differences with
just five seeds.

\paragraph{Reading results honestly.} Three habits:
(i) report \emph{success rate and collisions beside every mean} --- a cell
that reaches 60\% of the time with a beautiful mean clearance is a failing
cell; (ii) compare cells that differ \emph{only} in the factor your question
is about (``ceteris paribus'' --- all E1 algorithms consume the same scan
and gains); (iii) long-format CSVs (one row per run) are the safe storage,
because any summary is then a trivial \code{GROUP BY}.

\paragraph{Go deeper.} Art Owen's in-progress Monte Carlo book
\cite{owen2019} is free online; Chapter~1 motivates exactly our setup
(estimate an expectation by sampling, report the error) and Chapter~9 covers
variance reduction --- of which paired seeds are the simplest instance. For
a 10-minute video, Steve Brunton's lecture on Monte Carlo sampling
\cite{bruntonmc} shows the estimator and its error bars; the same channel's
``Crude Monte Carlo'' demonstration \cite{bruntoncrude} is a worked example
of seeding and error shrinkage with sample size.

\subsection{ToF vs LiDAR: what the sensor choice does to a planner}
\label{sec:concept-tof}

E3 sweeps beam counts but the deeper hardware question is \emph{which
sensor technology}, and E6 (\S\ref{sec:e6}) answers it directly with a
matched design. The physics of the two contenders:

\begin{itemize}[leftmargin=1.6em]
\item \textbf{LiDAR} (e.g.\ a rotating 288-beam unit): time-of-flight laser
ranging over a wide FoV at high angular resolution. Accurate and dense
--- that is the ``$288$ beams, $\sigma_0 = 0.02$'' reference row --- but
mechanically complex and comparatively expensive.
\item \textbf{ToF arrays} (e.g.\ ST VL53-series): solid-state flash LiDAR
on a chip --- a fixed grid of zones (8--64) over a narrow FoV, no moving
parts, a few dollars, but coarse angular resolution, narrow field, and
range error that \emph{grows with distance} (multipath and return-strength
effects). That growth is what the \code{'tof'} noise model encodes:
$s_b = \sigma_0 + k_{\mathrm{rng}}\, d_b$ inside a $60^\circ$ arc that is
fixed to the \emph{body} (the fan turns with the payload heading).
\end{itemize}

Two experimental-design warnings before reading E6's tables. First, the
comparison must be \emph{matched}: same world, same seeds, same algorithm,
same $\sigma_0$ --- only the sensor kind differs --- otherwise the effect of
``sensor'' is confounded with everything else. (This is
\S\ref{sec:gridconcept}'s rule~1 applied a second time.) Second, distinguish
the two axes of ``less sensor'': \emph{resolution} (beam count, E3) and
\emph{noise structure} (homoscedastic vs range-growing, E6). A sensor can
lose beams without losing reliability --- the $p_{\mathrm{unknown}}$
floor keeps unobserved bearings from being read as confidently free --- but
range-growing noise is a different attack: it blurs \emph{far} obstacles,
precisely where early warning lives. When you read E6's numbers, keep the
two effects a sensor change mixes apart: \emph{angular coverage} (FoV, beam
count) and \emph{per-return noise growth} ($s_b$). They can pull in opposite
directions: a narrow fan sees fewer side obstacles to react to, but it is
also blind to anything approaching from the side or rear --- which is where
the mover-contact column of Table~\ref{tab:e6} becomes the one to watch.

\paragraph{Go deeper.} ST's VL53L5CX product page \cite{stvl53} documents a
typical 8$\times$8 ToF array (FoV, zones, range-accuracy curve --- compare
with \code{sensor\_noise.m}'s parameters); SEEED's TF-Luna wiki \cite{tfluna}
is the cheapest reverse example (single-beam ToF, serial output). For the
planner's side of the trade, re-read \S\ref{sec:concept-unc} with ``which
$s_b$ curve am I feeding it?'' in mind.

\subsection{Self-check: five questions you can now answer}

\begin{enumerate}[leftmargin=1.9em]
\item Why can \code{apf} fail \emph{silently} (no gate rejections, no goal)
while a noisy deterministic histogram fails \emph{loudly} (many rejections)? \emph{(Local minimum vs
thresholded noise: \S\ref{sec:concept-apf}, \S\ref{sec:concept-unc}.)}
\item Which equation guarantees P-VFH reduces to VFH for a perfect sensor,
and which figure shows that as data? \emph{(Eq.~\eqref{eq:ukernel-concept};
Fig.~\ref{fig:paired} at low $\sigma_0$.)}
\item Your ToF array drops from 16 to 8 beams. What keeps the bearings no
beam covers from being read as confidently free, and which experiment
measures the cost? \emph{(The $p_{\mathrm{unknown}}$ floor; E3 and E6 ---
watch the mover-contact column.)}
\item Two algorithms share seed 3. Their min-clearance difference on that
seed is 0.4~m. Can noise explain it? \emph{(No --- paired design,
\S\ref{sec:concept-mc}.)}
\item Why is a ``flat goal rate, sagging clearance'' pattern a warning light
rather than good news? \emph{(The margin is what absorbs the next surprise;
\S\ref{sec:e3} how-to-read.)}
\end{enumerate}

\paragraph{References.} The complete reference list for the whole report
(videos included) is printed once at the end of the document.
